\documentclass[12pt]{article}

\usepackage[utf8]{inputenc}
\usepackage[spanish]{babel}
\usepackage[shortlabels]{enumitem}

\usepackage{mathtools}
\usepackage{amssymb}
\usepackage{amsthm}

\usepackage{geometry}
\geometry{margin=2.5cm}

\usepackage{hyperref}

\title{Solución Examen Final de Álgebra}
\author{}

\begin{document}

\section*{Solución Problema 1}

Consideramos el endomorfismo
\[
f:P_2(x)\longrightarrow P_2(x)
\]
tal que
\[
\ker f=\{a+bx+cx^2\in P_2(x):a=b=c\}
\]
y los polinomios con término independiente nulo se transforman en sí mismos.

La base canónica de \(P_2(x)\) es
\[
B_C=\{1,x,x^2\}.
\]

Como
\[
\ker f=\{a(1+x+x^2):a\in\mathbb{R}\}
=L<\{1+x+x^2\}>,
\]
tenemos
\[
f(1+x+x^2)=0.
\]

Además, los polinomios con término independiente nulo son de la forma
\[
bx+cx^2,
\]
y por hipótesis se transforman en sí mismos. Por tanto,
\[
f(x)=x,
\qquad
f(x^2)=x^2.
\]

Como
\[
1=(1+x+x^2)-x-x^2,
\]
aplicando \(f\),
\[
f(1)=f(1+x+x^2)-f(x)-f(x^2)=0-x-x^2=-x-x^2.
\]

\subsubsection*{a) Matriz del endomorfismo en la base canónica}

Las columnas de la matriz \(M_{B_C}(f)\) son las coordenadas de
\(f(1)\), \(f(x)\) y \(f(x^2)\) en la base canónica.

Tenemos
\[
f(1)=-x-x^2,
\qquad
f(x)=x,
\qquad
f(x^2)=x^2.
\]

Luego
\[
[f(1)]_{B_C}=
\begin{pmatrix}
0\\-1\\-1
\end{pmatrix},
\qquad
[f(x)]_{B_C}=
\begin{pmatrix}
0\\1\\0
\end{pmatrix},
\qquad
[f(x^2)]_{B_C}=
\begin{pmatrix}
0\\0\\1
\end{pmatrix}.
\]

Por tanto,
\[
\boxed{
M_{B_C}(f)=
\begin{pmatrix}
0&0&0\\
-1&1&0\\
-1&0&1
\end{pmatrix}.
}
\]

Equivalentemente, si
\[
p(x)=a+bx+cx^2,
\]
entonces
\[
p(x)=a(1+x+x^2)+(b-a)x+(c-a)x^2,
\]
y por tanto
\[
f(p(x))=(b-a)x+(c-a)x^2.
\]

\subsubsection*{b) Expresión implícita de \(\operatorname{Im}f\)}

De la expresión anterior,
\[
f(a+bx+cx^2)=(b-a)x+(c-a)x^2.
\]

Por tanto, todo polinomio de la imagen tiene término independiente nulo.
Además, dados \(\alpha,\beta\in\mathbb{R}\), tomando por ejemplo
\[
a=0,\qquad b=\alpha,\qquad c=\beta,
\]
se obtiene
\[
f(\alpha x+\beta x^2)=\alpha x+\beta x^2.
\]

Luego
\[
\operatorname{Im}f=L<\{x,x^2\}>.
\]

En coordenadas respecto de la base canónica,
\[
\boxed{
\operatorname{Im}f=
\{(u,v,w)\in\mathbb{R}^3:u=0\}.
}
\]

Por tanto,
\[
\boxed{\dim(\operatorname{Im}f)=2.}
\]

Lo podemos razonar de dos formas:

\[
\operatorname{Im}f=L<\{x,x^2\}>,
\]
y los polinomios \(x\) y \(x^2\) son linealmente independientes, luego
\[
\dim(\operatorname{Im}f)=2.
\]

Por otro lado, por el teorema de la dimensión,
\[
\dim P_2(x)=\dim(\ker f)+\dim(\operatorname{Im}f).
\]

Como
\[
\dim P_2(x)=3
\qquad\text{y}\qquad
\dim(\ker f)=1,
\]
entonces
\[
\dim(\operatorname{Im}f)=3-1=2.
\]

\subsubsection*{c) Cálculo de \(f(S)\)}

Sea
\[
S=\{(\lambda-2\mu+\delta,\lambda-2\mu-\delta,\lambda-2\mu):
\lambda,\mu,\delta\in\mathbb{R}\}.
\]

Tomamos
\[
r=\lambda-2\mu.
\]
Entonces
\[
S=\{(r+\delta,r-\delta,r):r,\delta\in\mathbb{R}\}.
\]

Si interpretamos estos vectores como coordenadas en la base canónica,
\[
p(x)=(r+\delta)+(r-\delta)x+rx^2.
\]

Aplicando la expresión de \(f\),
\[
f(p(x))=((r-\delta)-(r+\delta))x+(r-(r+\delta))x^2.
\]

Luego
\[
f(p(x))=-2\delta x-\delta x^2.
\]

Por tanto,
\[
f(S)=\{(0,-2\delta,-\delta):\delta\in\mathbb{R}\}
=L<\{(0,2,1)\}>.
\]

En forma implícita,
\[
\boxed{
f(S)=\{(u,v,w)\in\mathbb{R}^3:u=0,\ v-2w=0\}.
}
\]

Así,
\[
\boxed{\dim f(S)=1.}
\]

Esta dimensión tiene sentido. En efecto,
\[
S=\{(r+\delta,r-\delta,r):r,\delta\in\mathbb{R}\}
\]
tiene dimensión \(2\). Además,
\[
\ker f=L<\{(1,1,1)\}>.
\]

El vector \((1,1,1)\) pertenece a \(S\), ya que se obtiene tomando
\[
r=1,
\qquad
\delta=0.
\]

Por tanto,
\[
\dim(S\cap\ker f)=1.
\]

Aplicando el teorema de la dimensión a la restricción
\[
f_{\mid S}:S\longrightarrow f(S),
\]
tenemos
\[
\dim f(S)=\dim S-\dim(S\cap\ker f)=2-1=1.
\]

\subsubsection*{d) Matriz de \(f\) en la base \(B^*\)}

Sea
\[
B^*=\{1+x+x^2,\ 1+x,\ 1\}.
\]

Calculamos las imágenes de los vectores de \(B^*\):
\[
f(1+x+x^2)=0,
\]
\[
f(1+x)=f(1)+f(x)=(-x-x^2)+x=-x^2,
\]
\[
f(1)=-x-x^2.
\]

Ahora expresamos cada imagen en coordenadas respecto de \(B^*\).

Para \(f(1+x+x^2)=0\),
\[
[f(1+x+x^2)]_{B^*}=
\begin{pmatrix}
0\\0\\0
\end{pmatrix}.
\]

Para \(-x^2\), buscamos \(\alpha,\beta,\gamma\in\mathbb{R}\) tales que
\[
-x^2=\alpha(1+x+x^2)+\beta(1+x)+\gamma.
\]

Comparando coeficientes:
\[
\begin{cases}
\alpha+\beta+\gamma=0,\\
\alpha+\beta=0,\\
\alpha=-1.
\end{cases}
\]

De aquí,
\[
\alpha=-1,
\qquad
\beta=1,
\qquad
\gamma=0.
\]

Luego
\[
[-x^2]_{B^*}=
\begin{pmatrix}
-1\\1\\0
\end{pmatrix}.
\]

Para \(-x-x^2\), buscamos
\[
-x-x^2=\alpha(1+x+x^2)+\beta(1+x)+\gamma.
\]

Comparando coeficientes:
\[
\begin{cases}
\alpha+\beta+\gamma=0,\\
\alpha+\beta=-1,\\
\alpha=-1.
\end{cases}
\]

De aquí,
\[
\alpha=-1,
\qquad
\beta=0,
\qquad
\gamma=1.
\]

Luego
\[
[-x-x^2]_{B^*}=
\begin{pmatrix}
-1\\0\\1
\end{pmatrix}.
\]

Así,
\[
\boxed{
M_{B^*}(f)=
\begin{pmatrix}
0&-1&-1\\
0&1&0\\
0&0&1
\end{pmatrix}.
}
\]

También se puede obtener matricialmente mediante cambio de base.
La matriz de cambio de coordenadas de \(B^*\) a \(B_C\) es
\[
P=
\begin{pmatrix}
1&1&1\\
1&1&0\\
1&0&0
\end{pmatrix},
\]
ya que sus columnas son las coordenadas de los vectores de \(B^*\) en la base canónica.

Entonces
\[
M_{B^*}(f)=P^{-1}M_{B_C}(f)P.
\]

Calculando,
\[
P^{-1}M_{B_C}(f)P=
\begin{pmatrix}
0&-1&-1\\
0&1&0\\
0&0&1
\end{pmatrix}.
\]

\section*{Solución Problema 2}

Sea
\[
A=
\begin{pmatrix}
4&-1&6\\
2&1&6\\
2&-1&8
\end{pmatrix}.
\]

\subsubsection*{a) Estudio de diagonalización}

Calculamos el polinomio característico:
\[
p_A(\lambda)=\det(\lambda I-A).
\]

Tenemos
\[
\lambda I-A=
\begin{pmatrix}
\lambda-4&1&-6\\
-2&\lambda-1&-6\\
-2&1&\lambda-8
\end{pmatrix}.
\]

Por tanto,
\[
p_A(\lambda)=\lambda^3-13\lambda^2+40\lambda-36.
\]

Factorizamos:
\[
\lambda^3-13\lambda^2+40\lambda-36
=(\lambda-2)^2(\lambda-9).
\]

Luego los autovalores son
\[
\lambda_1=2
\quad\text{con multiplicidad algebraica }2,
\]
y
\[
\lambda_2=9
\quad\text{con multiplicidad algebraica }1.
\]

Estudiamos los subespacios propios.

Para \(\lambda=2\),
\[
A-2I=
\begin{pmatrix}
2&-1&6\\
2&-1&6\\
2&-1&6
\end{pmatrix}.
\]

El sistema \((A-2I)X=0\) queda
\[
2x-y+6z=0.
\]

Por tanto,
\[
y=2x+6z.
\]

Así,
\[
E_2=\{(x,2x+6z,z):x,z\in\mathbb{R}\}.
\]

Separando parámetros:
\[
(x,2x+6z,z)=x(1,2,0)+z(0,6,1).
\]

Luego
\[
E_2=L<\{(1,2,0),(0,6,1)\}>,
\]
y
\[
\dim E_2=2.
\]

Para \(\lambda=9\),
\[
A-9I=
\begin{pmatrix}
-5&-1&6\\
2&-8&6\\
2&-1&-1
\end{pmatrix}.
\]

Resolviendo \((A-9I)X=0\), se obtiene
\[
x=y=z.
\]

Por tanto,
\[
E_9=L<\{(1,1,1)\}>,
\]
y
\[
\dim E_9=1.
\]

Como
\[
\dim E_2+\dim E_9=2+1=3,
\]
existe una base de \(\mathbb{R}^3\) formada por autovectores de \(A\).

Por tanto,
\[
\boxed{A\text{ es diagonalizable}.}
\]

\subsubsection*{b) Matrices \(P\) y \(D\)}

Tomamos como columnas de \(P\) una base de autovectores:
\[
v_1=(1,2,0),
\qquad
v_2=(0,6,1),
\qquad
v_3=(1,1,1).
\]

Los dos primeros corresponden al autovalor \(2\), y el tercero corresponde al autovalor \(9\).

Por tanto,
\[
P=
\begin{pmatrix}
1&0&1\\
2&6&1\\
0&1&1
\end{pmatrix}.
\]

La matriz diagonal asociada es
\[
D=
\begin{pmatrix}
2&0&0\\
0&2&0\\
0&0&9
\end{pmatrix}.
\]

Con esta elección,
\[
\boxed{
P^{-1}AP=D.
}
\]

\subsubsection*{c) Interpretación del proceso de diagonalización}

Diagonalizar \(A\) consiste en cambiar de la base canónica a una base formada por autovectores del endomorfismo.

En la base canónica, el endomorfismo está representado por
\[
A=
\begin{pmatrix}
4&-1&6\\
2&1&6\\
2&-1&8
\end{pmatrix}.
\]

Sin embargo, en la base
\[
B=\{(1,2,0),(0,6,1),(1,1,1)\},
\]
el endomorfismo actúa de forma mucho más simple:
\[
f(1,2,0)=2(1,2,0),
\]
\[
f(0,6,1)=2(0,6,1),
\]
\[
f(1,1,1)=9(1,1,1).
\]

Por eso, en esa base, la matriz asociada es diagonal:
\[
D=
\begin{pmatrix}
2&0&0\\
0&2&0\\
0&0&9
\end{pmatrix}.
\]

La matriz \(P\) es la matriz de cambio de base: sus columnas son los vectores de la nueva base escritos en la base canónica.

\section*{Solución Cuestión 1}

Sea
\[
M=
\begin{pmatrix}
1&-1\\
-1&1
\end{pmatrix}
\]
y sea
\[
X=
\begin{pmatrix}
a&b\\
c&d
\end{pmatrix}.
\]

Entonces
\[
MX=
\begin{pmatrix}
1&-1\\
-1&1
\end{pmatrix}
\begin{pmatrix}
a&b\\
c&d
\end{pmatrix}
=
\begin{pmatrix}
a-c&b-d\\
-a+c&-b+d
\end{pmatrix}.
\]

Por tanto,
\[
\operatorname{traza}(MX)=(a-c)+(-b+d)=a-b-c+d.
\]

La condición que define \(U\) es
\[
a-b-c+d=0.
\]

Así,
\[
d=-a+b+c.
\]

Luego
\[
X=
\begin{pmatrix}
a&b\\
c&-a+b+c
\end{pmatrix}.
\]

Separando parámetros:
\[
X=
a
\begin{pmatrix}
1&0\\
0&-1
\end{pmatrix}
+b
\begin{pmatrix}
0&1\\
0&1
\end{pmatrix}
+c
\begin{pmatrix}
0&0\\
1&1
\end{pmatrix}.
\]

Por tanto, una base de \(U\) es
\[
\boxed{
\left\{
\begin{pmatrix}
1&0\\
0&-1
\end{pmatrix},
\begin{pmatrix}
0&1\\
0&1
\end{pmatrix},
\begin{pmatrix}
0&0\\
1&1
\end{pmatrix}
\right\}.
}
\]

Además,
\[
\boxed{\dim U=3.}
\]

\section*{Solución Cuestión 2}

La aplicación lineal
\[
f:\mathbb{R}^2\longrightarrow\mathbb{R}^3
\]
viene dada por
\[
f(1,0)=(2,1,2),
\qquad
f(0,1)=(3,2,1).
\]

Luego su matriz en las bases canónicas es
\[
M(f)=
\begin{pmatrix}
2&3\\
1&2\\
2&1
\end{pmatrix}.
\]

La aplicación
\[
g:\mathbb{R}^3\longrightarrow\mathbb{R}^2
\]
viene dada por
\[
g(1,0,0)=(3,-1),
\]
\[
g(0,1,0)=(1,0),
\]
\[
g(0,0,1)=(1,-2).
\]

Luego
\[
M(g)=
\begin{pmatrix}
3&1&1\\
-1&0&-2
\end{pmatrix}.
\]

Como
\[
h=g\circ f,
\]
entonces
\[
M(h)=M(g)M(f).
\]

Calculamos:
\[
M(h)=
\begin{pmatrix}
3&1&1\\
-1&0&-2
\end{pmatrix}
\begin{pmatrix}
2&3\\
1&2\\
2&1
\end{pmatrix}
=
\begin{pmatrix}
9&12\\
-6&-5
\end{pmatrix}.
\]

Por tanto, matricialmente:
\[
h(4,-2)=
\begin{pmatrix}
9&12\\
-6&-5
\end{pmatrix}
\begin{pmatrix}
4\\
-2
\end{pmatrix}
=
\begin{pmatrix}
12\\
-14
\end{pmatrix}.
\]

Luego
\[
\boxed{h(4,-2)=(12,-14).}
\]

Por coordenadas:
\[
f(4,-2)=4f(1,0)-2f(0,1).
\]

Así,
\[
f(4,-2)=4(2,1,2)-2(3,2,1)=(8,4,8)-(6,4,2)=(2,0,6).
\]

Entonces
\[
h(4,-2)=g(f(4,-2))=g(2,0,6).
\]

Como \(g\) es lineal,
\[
g(2,0,6)=2g(1,0,0)+6g(0,0,1).
\]

Por tanto,
\[
g(2,0,6)=2(3,-1)+6(1,-2)=(6,-2)+(6,-12)=(12,-14).
\]

\section*{Solución Cuestión 3}

Tenemos
\[
U=L<\{(1,0,0,1),(1,1,1,1),(0,2,2,0)\}>
\]
y
\[
W=\{(x,y,z,t)\in\mathbb{R}^4:
3x+y-z-3t=0,\ y-t=0\}.
\]

Observamos que
\[
(0,2,2,0)=2(1,1,1,1)-2(1,0,0,1).
\]

Por tanto,
\[
U=L<\{(1,0,0,1),(1,1,1,1)\}>.
\]

Un vector de \(U\) es de la forma
\[
\alpha(1,0,0,1)+\beta(1,1,1,1)
=(\alpha+\beta,\beta,\beta,\alpha+\beta).
\]

Así,
\[
U=\{(x,y,z,t)\in\mathbb{R}^4:x=t,\ y=z\}.
\]

Para que \(U+W\) sea suma directa, debe cumplirse
\[
U\cap W=\{(0,0,0,0)\}.
\]

Pero
\[
(1,1,1,1)\in U.
\]

Además,
\[
3\cdot 1+1-1-3\cdot 1=0
\]
y
\[
1-1=0.
\]

Luego
\[
(1,1,1,1)\in W.
\]

Por tanto,
\[
(1,1,1,1)\in U\cap W
\]
y este vector no es el vector nulo.

Así,
\[
U\cap W\neq \{(0,0,0,0)\}.
\]

Concluimos que
\[
\boxed{U+W\text{ no es suma directa}.}
\]

\section*{Solución Cuestión 4}

\subsubsection*{1. Una matriz \(A\) tiene autovalor \(0\) si y solo si \(\det(A)=0\)}

La afirmación es verdadera.

En efecto, \(0\) es autovalor de \(A\) si y solo si existe un vector no nulo \(v\) tal que
\[
Av=0v=0.
\]

Esto equivale a que el sistema homogéneo
\[
Av=0
\]
tenga soluciones no nulas, es decir, a que \(A\) no sea invertible.

Una matriz cuadrada no es invertible si y solo si
\[
\det(A)=0.
\]

Por tanto,
\[
\boxed{\text{la afirmación es verdadera}.}
\]

\subsubsection*{2. Semejanza de matrices}

Sean
\[
A=
\begin{pmatrix}
1&2&0\\
-1&3&1\\
0&1&1
\end{pmatrix}
\]
y
\[
B=
\begin{pmatrix}
5&0&-4\\
0&3&0\\
2&0&-1
\end{pmatrix}.
\]

Si dos matrices son semejantes, entonces tienen la misma traza.

Calculamos:
\[
\operatorname{traza}(A)=1+3+1=5.
\]

Por otro lado,
\[
\operatorname{traza}(B)=5+3+(-1)=7.
\]

Como
\[
\operatorname{traza}(A)\neq \operatorname{traza}(B),
\]
las matrices no pueden ser semejantes.

Por tanto,
\[
\boxed{\text{la afirmación es falsa}.}
\]

\end{document}
